% sec4b.tex — c05｜原文 p25（Corollary 4.13 的证明起）–p30（References 之前）｜§4.3–§4.5
% 公式号 (51)–(56)；编号环境 4.14–4.21
% 块边界：本块从 p25 上 Corollary 4.13 的 Proof 起。该推论的陈述、注 4.12，以及 p25 上半
%   命题 4.11 的证明，均在 sec4a.tex（c04）；c04 负责 (39)–(50)。
% 注 4.14 与式 (52) 位于 p26 页首（按页面位置归本块；BLOCK_TAGS 亦把 (52) 记在 sec4b）。
% 图 7 在原书浮动于 p28 顶部，图 8 浮动于 p29 顶部：两幅均不复刻。

% [原文 p.25]
\begin{proof}
第一个断言由把例 4.2 应用于命题 4.11 得到。最后一个断言由取手征理论所对应的典范超对称理论得到（见例 4.2），前提是每个整模形式都作为某个手征理论的配分函数出现。为验证这一遗留断言，设 $f(q)\in\mathrm{MF}^{\Z}$ 是权为 $n/2$ 的模形式。若 $n$ 为奇数，则 $f(q)=0$，可取零理论作为我们的手征理论。若 $n$ 为偶数，则考虑 $f(q)\eta^n=q^{n/24}\sum a_nq^n\in q^{n/24}\Z((q))$。定义
\begin{equation}
\cH(k)=(\mathrm{Cl}_2\C^{1|1})^{\otimes n/2}\otimes\C^{\times a_n},
\qquad
H^{(r)}_+=\bigoplus_{k=-1}^{\infty}\frac{1}{r}\left(k+\frac{n}{24}\right)\cdot\id_{\cH(k)}
\tag{51}
\end{equation}
其中 $\mathrm{Cl}_2(\C^{1|1})_{\C}$ 是实现 $\mathrm{Cl}_2$ 与 $\C$ 之间 Morita 等价的双模。这个双模的 Clifford 超迹等于 1，故（利用张量积之迹的乘性）理论 (51) 的配分函数为
\[
\begin{aligned}
Z(r,\tau)
&:=\Tr_{\mathrm{Cl}_n}\left(q^{H^{(r)}_+}\bar q^{H^{(r)}_-}\right)
=\Tr_{\mathrm{Cl}_n}\left(q^{H^{(r)}_+}\right)\\
&=\sum_{k\in\Z}\Tr_{\mathrm{Cl}_n}\left((\mathrm{Cl}_2\C^{1|1})^{\otimes n/2}\right)\cdot\Tr\left(q^{k+n/24}\cdot\id_{\C^{\times a_n}}\right)
=q^{n/24}\sum_{k\in\Z}a_kq^k,
\end{aligned}
\]
它在 $\MP_2(\Z)$-作用下按 $n$ 次单位根变换，这里用到 $f(q)$ 的模性与注 4.1。
\end{proof}
% [Tutor Note: 式 (51) 右式的下指标原文印作 $k=-1$，但 $\cH(k)$ 的定义与 $k$ 无关，且同一段的分拆与下面的配分函数求和都遍历 $k\in\Z$；此下指标疑为 $k\in\Z$（或 $k=-\infty$）之误，照排不改，已登记 ERRATA。]
% [Tutor Note: 式 (51) 中的 $\C^{\times a_n}$ 表示 $a_n$ 个 $\C$ 的直和，故 $\Tr(q^{k+n/24}\id_{\C^{\times a_n}})=a_n$，这与下一行给出的系数 $a_k$ 相符，可用来核对记号。]
% [Tutor Note: 式 (51) 里花体 $\cH(k)$ 与斜体 $H^{(r)}_+$ 是两个不同对象：前者是权空间，后者是圆周长 $2\pi r$ 上的哈密顿量，$\frac1r(k+\frac n{24})$ 即第 $k$ 个 Fourier 模的能级加 Casimir 项。]
% [Tutor Note: 式 (51) 中印作 $(\mathrm{Cl}_2\C^{1|1})^{\otimes n/2}$，紧接的正文却把同一双模记作 $\mathrm{Cl}_2(\C^{1|1})_{\C}$（左 $\mathrm{Cl}_2$、右 $\C$）；两种写法指同一个 Morita 双模。]

% [原文 p.26]
% 注 4.14 全文位于本页页首（含式 (52)），紧接上一页的推论 4.13 证明，归本块。
% [Tutor Note: 注 4.14 引 [GPPV21, Table 2]，但本综述全文没有实体表格（既无 Table 1 也无 Table 2），亦无表格环境；照译该引用、不补表，读者应到 [GPPV21] 原文中查表。]
\begin{remark}[4.14]
推论 4.13 敲响了警钟：边缘同态
\begin{equation}
\pi_*\TMF\to\mathrm{MF}^{\Z}
\tag{52}
\end{equation}
从 TMF 的系数到整模形式，不是满射 [Hop02, Proposition 4.6]。因此，若猜想 1.1 要成立，就不应把全部整模形式都实现为超对称场论的配分函数。特别地，若采用定义 4.9，猜想 1.1 是假的。值得注意的是，物理中最确凿的那些超对称场论例子都与 (52) 的像一致，例如见 [GPPV21, Table 2]。然而，把给定的整模形式实现为一段场论故事这一问题，在物理文献中尚未完全解决，见 [LP23]。数学上的预期是：某些（手征）理论无法按下文 §4.5 的意义向下扩张，而向下扩张未必唯一。这既使某些模形式无法实现为配分函数，也使给定的模形式可以有多种方式实现为某个场论的配分函数。这一预期与 (52) 既不满射也不单射相符。
\end{remark}
% [Tutor Note: 注 4.14 与推论 4.13 并不矛盾：4.13 只断言每个整模形式可由定义 4.9 的（手征）理论实现；4.14 指出定义 4.9 若当作猜想 1.1 的场论，就要求 $\pi_*\TMF\to\mathrm{MF}^{\Z}$ 满射，而它不满射；弥合这一落差正是 §4.5 的向下扩张。]
% [Tutor Note: 「边缘同态」在这里指谱的系数到同伦群/模形式的边缘映射，与 §2–§3 中把上同调理论送到系数的边缘映射同义；它不满射是 TMF 与整模形式之间真实的差异，不是定义 4.9 的记号问题。]

\subsection{重新包装为配边范畴的表示}
可以把定义 4.6 与 4.9 中的分析数据重新包装，分别作为配边范畴 $\EBord_2$ 与 $\EBord_{2|1}$ 的表示的一部分。下面简要回顾这些配边范畴，细节参见 [ST11]。

$\EBord_2$ 的对象是带度量的圆周 $S^R_r$ 与 $S^{\mathrm{NS}}_r$ 的有限不交并，并带 spin 结构；其中 $r\in\R_{>0}$ 记录半径，$R$ 代表 Ramond（别名：周期的、奇的、或非边界的）spin 结构，$\mathrm{NS}$ 代表 Neveu–Schwarz（别名：反周期的、偶的、或边界的）spin 结构。$\EBord_2$ 的态射是带 spin 结构与测地边界的平坦 2-流形，如图 7 所示。除图 5 中的柱面与环面之外，改变柱面的源与目标数据就给出图 7 中的「通心粉」。由 §4.1 开头的讨论，固定源与目标的欧几里得通心粉的模空间是 $\HH/r\Z$，其中 $r\in\R_{>0}$ 记录半径，$\tau=\frac{x+it}{2\pi r}\in\HH$ 给出柱面的高度 $t$，而 $x$ 是边界参数化之差。复合对应于 $\HH$ 中的加法，例如 $C^R_{r,\tau}\circ C^R_{r,\tau'}\simeq C^R_{r,\tau+\tau'}$。环面 $\C/2\pi r\Z\oplus2\pi\tau\Z$ 上有 4 个 spin 结构。metaplectic 群 $\MP_2(\Z)$ 的作用置换三个偶 spin 结构，而奇 spin 结构在此作用下保持不变。我们强调，平坦度量排除了亏格更高的闭曲面。平坦度量连同测地边界排除了不同尺寸、或不同条数的圆周之间的配边，例如裤子不是 $\EBord_2$ 中的态射。
% [Tutor Note: $\MP_2(\Z)$ 是 $\mathrm{SL}_2(\Z)$ 的二重覆盖（metaplectic 群）；权为半整数的模形式在其上变换，而 (52) 中的整模形式权为整数、只在 $\mathrm{SL}_2(\Z)$ 上变换。本条与证明末「按 $n$ 次单位根变换」相呼应。]
% [Tutor Note: 环面上的 4 个 spin 结构按「偶/奇」分为 3+1：三个偶（可边界）结构被 $\MP_2(\Z)$ 置换，奇结构固定。这里的奇/偶与 $R$/$\mathrm{NS}$ 的别名「奇/偶」一致，勿与 $\Z/2$-分次混同。]

加上超对称，$\EBord_{2|1}$ 的对象是 $1|1$ 维超圆周；其模空间同样由 spin 结构与半径参数化 [BE23c, §4.5]，故沿用同一记号 $S^R_r,S^{\mathrm{NS}}_r\in\EBord_{2|1}$。它们之间的配边是带测地边界的 $2|1$ 维欧几里得流形。与 $1|1$ 维情形类似，这在 $\EBord_{2|1}$ 的态射模空间中引入奇方向。例如，Ramond 扇区中的超欧几里得柱面由 $\R_{>0}\times\HH^{2|1}$ 参数化，其复合为
\begin{equation}
C^R_{r,\tau,\bar\tau,\theta}\circ C^R_{r,\tau',\bar\tau',\theta'}\simeq C^R_{r,\tau+\tau',\bar\tau+\bar\tau'+\theta\theta',\theta+\theta'},
\qquad r\in\R_{>0},\ (\tau,\bar\tau,\theta)\in\HH^{2|1}.
\tag{53}
\end{equation}
这给出超半群（supersemigroup）表示 (43) 的几何对应物。
% [Tutor Note: (53) 中 $\bar\tau$ 的平移量 $\theta\theta'$ 是 $\HH^{2|1}$ 半群乘法的奇-奇项，$\theta$ 的平移量则是 $\theta+\theta'$；把它与式 (43) 的代数侧逐项对照，可确认「柱面的复合」与「超半群乘法」是同一结构的两面。]

为建立与猜想 1.1 相关的次数 $n$ 超欧几里得场论，首先必须描述次数扭。以 $\mathrm{Cl}(L_r\R^n)$ 记向量空间 $L_r\R^n=C^\infty(S^1_r;\R^n)$ 关于双线性配对 $\int_0^{2\pi r}\langle\psi(t),\varphi(t)\rangle dt$ 的 Clifford 代数。$S^1_r$ 在 $L_r\R^n$ 上的旋转作用保持此双线性配对，从而给出 $\mathrm{Cl}(L_r\R^n)$ 上由代数自同构构成的 $S^1_r$-作用。对每个 $\tau\in\HH/r\Z$，存在代数映射 $\varphi_r(\tau)\colon\mathrm{Cl}(L_r\R^n)\to\mathrm{Cl}(L_r\R^n)$，它由该 $S^1_r$-作用的解析延拓得到。细节参见 [ST11, page 59]。

% [原文 p.27]
\begin{definition}[4.15（次数扭，仅 Ramond 扇区的略述）]
次数扭的取值是
\begin{equation}
\begin{gathered}
\alpha_n\colon\EBord_{2|1}\to\Alg,\\
\alpha_n(S^R_r)=\mathrm{Cl}(L_r\R^n),\\
\alpha_n(C^R_{r,\tau,\bar\tau,\theta})={}_{\varphi_r(\tau)\mathrm{Cl}(L_r\R^n)}\mathrm{Cl}(L_r\R^n)_{\mathrm{Cl}(L_r\R^n)},\\
\alpha_n\Big(\,\tikz[baseline=-0.35ex,line width=0.35pt,scale=0.32]{
\draw (0,0) ellipse (0.42 and 0.18);
\draw (2.2,0) ellipse (0.42 and 0.18);
\draw (0,0.14) .. controls (0.35,1.0) and (1.85,1.0) .. (2.2,0.14);
\draw (0.42,0.12) .. controls (0.75,0.75) and (1.45,0.75) .. (1.78,0.12);
\node[anchor=base west,inner sep=1.4pt] at (2.62,0.02) {$\scriptstyle r,\tau,\bar\tau,\theta$};
}\,\Big)=\mathrm{Cl}(L_r\R^n)_{\varphi_r(\tau)\mathrm{Cl}(L_r\R^n)\otimes\mathrm{Cl}(L_r\R^n)},\\
\alpha_n\Big(\,\tikz[baseline=-0.35ex,line width=0.35pt,scale=0.32]{
\draw (0,0) ellipse (0.42 and 0.18);
\draw (2.2,0) ellipse (0.42 and 0.18);
\draw (0,-0.14) .. controls (0.35,-1.0) and (1.85,-1.0) .. (2.2,-0.14);
\draw (0.42,-0.12) .. controls (0.75,-0.75) and (1.45,-0.75) .. (1.78,-0.12);
\node[anchor=base west,inner sep=1.4pt] at (2.62,-0.02) {$\scriptstyle r,\tau,\bar\tau,\theta$};
}\,\Big)={}_{\varphi_r(\tau)\mathrm{Cl}(L_r\R^n)\otimes\mathrm{Cl}(L_r\R^n)}\mathrm{Cl}(L_r\R^n),\\
\alpha_n\Big(\,\tikz[baseline=-0.35ex,line width=0.35pt,scale=0.32]{
\draw (0,0) circle (0.48);
\draw (0,0) circle (0.24);
\node[anchor=base west,inner sep=1.4pt] at (0.62,0) {$\scriptstyle r,\tau,\bar\tau,\theta$};
}\,\Big)=\mathrm{Pf}^{\otimes-n}
\end{gathered}
\tag{54}
\end{equation}
其中所有模都来自 $\mathrm{Cl}(L_r\R^n)$ 作为其自身的双模，而左作用被代数自同构 $\varphi_r(\tau)$ 扭过；$\mathrm{Pf}$ 是 Pfaffian 线丛，见注 4.1。这些取值与 $(\bar\tau,\theta)$ 无关，因此次数扭是手征的。
\end{definition}
% [Tutor Note: (54) 第 3、4、5 行的自变量，原文印的是三个手绘小配边图（无文字符号）：起拱的把手（两个边界圆 + 上拱的带）、其镜像（下凹的带）、同心双圆。译本以同构形的简笔线图复刻，只为可读性，不含原文之外的任何信息；下标 $r,\tau,\bar\tau,\theta$ 的位置照原文紧跟在图右侧。]
% [Tutor Note: (54) 用的是双模记号 ${}_AM_B$：左下标记左作用的扭、右下标记右作用的扭，缺省即该侧未扭。第 2 行读作 ${}_{\varphi_r(\tau)\mathrm{Cl}}\mathrm{Cl}_{\mathrm{Cl}}$；第 3 行（杯形把手）缺左下标、第 4 行（帽形把手）缺右下标，二者恰好把左扭与右扭对调，不可互换。]

\begin{remark}[4.16]
「超化」（superfication）函子 $\EBord_2\to\EBord_{2|1}$ [ST11, Eq. 4.14] 使人能把扭 $\alpha\colon\EBord_{2|1}\to\Alg$ 限制为扭 $\EBord_2\to\Alg$。特别地，扭 (54) 决定了 2 维欧几里得配边范畴的扭。
\end{remark}

\begin{remark}[4.17]
定义 4.15 省略了 Neveu–Schwarz 的取值，但它们来自向量空间 $C^\infty(S^1;\mu\otimes\R^n)$ 的类似 Clifford 代数，其中 $\mu$ 是 Möbius 丛，即 Neveu–Schwarz spin 结构对应的旋量丛。作为从欧几里得配边范畴出发的函子，次数扭的详细略述见 [ST11, §4.3]，另见 [Lud22]。该理论是手征的，因此预期容许一个典范的超对称扩张，推广例 4.2。然而，这一扩张在 [ST11] 的配边框架中的细节尚未做出。
\end{remark}

令 $\QFT^{\alpha}_{2|1}(\cX)$ 记 $\cX$ 上 $\alpha$-扭的反射正 $2|1$ 维欧几里得场论构成的群胚（见定义 1.9）。在存在手征超对称时，反射正性更为微妙，但最终反射结构的数据遵循与定义 2.13 类似的模式。我们先分析点上的场论之群胚 $\QFT^n_2:=\QFT^{\alpha_n}_2(\mathrm{pt})$ 与 $\QFT^n_{2|1}:=\QFT^{\alpha_n}_{2|1}(\mathrm{pt})$，其中的场论被次数扭 (54) 所扭。

\begin{hypothesis}[4.18（{[ST04, Theorem 1.0.2] [Che08, Theorem 2.2.2]}）]
存在叠的映射
\begin{equation}
\QFT^n_2\to\qft^n
\qquad\text{与}\qquad
\QFT^n_{2|1}\to\sqft^n
\tag{55}
\end{equation}
它们由限制到 Ramond 扇区得到。连同定理 4.10 与 4.11，叠 $\QFT^n_{2|1}$ 容许到表示 $K((q))$ 的谱中第 $n$ 个空间的映射，以及到整模形式集合的映射。
\end{hypothesis}
% [Tutor Note: 假设 4.18 原文 "There are map of stacks" 有语法笔误，应为 "There are maps of stacks"；中文无单复数，照译，已登记 ERRATA。]
% [Tutor Note: (55) 的两个映射都是「限制到 Ramond 扇区」，即忘掉超对称与 Neveu–Schwarz 扇区的数据；这正是注 4.14 之后所说的反例来源——配边定义严格比定义 4.9 携带更多信息（见 §4.4 开头）。]

\begin{proof}[略述]
我们聚焦于 $\QFT^n_{2|1}\to\sqft^n$；无超对称的情形更容易。

首先证明：限制到 Ramond 扇区的 (54) 的扭场论决定了通常 Clifford 代数 $\mathrm{Cl}_n$ 上的模，做法依照 [ST11, §6]。想法是按 Fourier 模，即 $S^1_r$-作用的权空间，分解代数 $\mathrm{Cl}(L_r\R^n)$。这些权给出稠密包含
\[
\mathrm{Cl}_n\otimes\bigotimes_{k>0}\mathrm{Cl}(\C^n_k\oplus\C^n_{-k})\hookrightarrow\mathrm{Cl}(L_r\R^n),
\]
其中 $\mathrm{Cl}_n$ 是 $S^1_r$-不动点（即常值映射）的 Clifford 代数，而 $\mathrm{Cl}(\C^n_k\oplus\C^n_{-k})$ 是 Fourier 数为 $k$ 与 $-k$ 的环路构成的 Clifford 代数。对配对 $\int_{S^1}\langle\varphi(t),\psi(t)\rangle dt$，这样的环路只彼此非平凡配对，且该配对是双曲配对。于是得到 Clifford 代数 $\mathrm{Cl}(\C^n_k\oplus\C^n_{-k})\subset\mathrm{Cl}(L_r\R^n)$ 作为子代数。此外，对每个 $k$，$\mathrm{Cl}(\C^n_k\oplus\C^n_{-k})\simeq\mathrm{Cl}_{n,n}$ 且 Morita 等价于 $\C$。Morita 双模的张量积的完备化给出从 $\mathrm{Cl}(L_r\R^n)$-模到 $\mathrm{Cl}_n$-模的所需函子，见 [ST11, Eq. 6.2]。

把这一点应用于 $E\in\QFT^n_{2|1}$ 在 Ramond 扇区超环带 (53) 的子范畴上的取值，得到形如 (43) 的 $\mathrm{Cl}_n$-线性超半群表示。这给出 $\sqft^n$ 中一个对象的数据。$E$ 在 Ramond–Ramond 超环面上的取值施加
% [原文 p.28]
模性条件 (49) [ST11, page 60]。应用与假设 2.14 类似的论证，扭场论之间的同构来自状态空间之间的射影同构。于是得到函子 $\QFT^n_{2|1}\to\sqft^n$。
\end{proof}
% [Tutor Note: 稠密包含里 $\mathrm{Cl}_n$ 来自常值环路（$k=0$）；配对只把模 $k$ 与 $-k$ 配起来，每个因子 $\mathrm{Cl}(\C^n_k\oplus\C^n_{-k})$ 上的配对非退化且双曲，故它 Morita 等价于 $\C$。]
% [Tutor Note: 原文 "with of Fourier numbers $k$ and $-k$" 多一个 of，属笔误；中文不受影响，已登记 ERRATA。]

% 原书中图 7 浮动于本页顶部，上一段的句子被它隔断；为避免割裂汉语句子，图 7 的
% \figomit 置于该段之后。图 7 的正文引用见 §4.3 首段「如图 7 所示」。
\figomit{7}{$2|1$ 维配边范畴的生成元与关系，类比于图 2 与图 4，但几何参数由 $(r,\tau,\bar\tau,\theta)\in\R_{>0}\times\HH^{2|1}$ 给出。在记录好源与目标之后，复合对应于半群 $\HH^{2|1}$ 中的乘法。按惯例，配边沿正（向上）方向读。左图：把一条「悲伤通心粉」粘到一条「快乐通心粉」上，得到一个柱面。右图：把柱面粘到一条「快乐通心粉」上，得到一条「更快乐的通心粉」。}
% [Tutor Note: 图 7 内还印出两条复合关系（对应 (53) 的复合律）：左为一端粘合给出 $C^R_{r,\tau+\tau',\bar\tau+\bar\tau'+\theta\theta',\theta+\theta'}$，右为两段柱面先竖向粘合再与通心粉复合；这些式子属于图的内容，随图一并省略。]

\subsection{配边框架的更多洞见}
与假设 3.20 相反，假设 4.18 中的映射 $\QFT^n_{2|1}\to\sqft^n$ 不再是叠的满射：配边定义所包含的信息严格多于定义 4.9，例如 $E(S^{\mathrm{NS}}_r)$ 在 Neveu–Schwarz 扇区上的取值。下面强调若干进一步的增强。

\begin{example}[4.19（实结构）]
在 [ST11] 的配边框架中，$\EBord_{2|1}(\cX)$ 的对象与态射都构成叠，意思是配边的等距群被编码在结构之中。这导致对象上一个在 1 维情形中不存在的额外 $\Z/2$-自同构，归根结底是因为 $-\id\in\mathrm{O}_2$ 是 $\R^2$ 的保向等距（与 $\R$ 上反转定向的 $-1\in\mathrm{O}(1)$ 形成对照）。物理上，这个自同构是宇称–时间反演。对反射正的理论，这个对称性决定状态空间与其复共轭之间的同构，即一个实结构的数据 [BE23c, §8.4]，从而把 4.18 中取值于 $K((q))$ 的上闭链映射细化为取值于 $\KO((q))$ 的映射。这一论证也可以用物理语言理解，见 [Wit87] 与 [GPPV21, 3.2.2]。
\end{example}
% [Tutor Note: 关键在行列式：$-\id\in\mathrm{O}_2$ 的行列式为 $(-1)^2=1$，故保向；而 $-1\in\mathrm{O}(1)$ 的行列式为 $-1$，反转定向。保向的额外自同构才允许它在反射正理论里给出状态空间到其复共轭的同构，从而把 $K((q))$ 细化为 $\KO((q))$。]

\begin{example}[4.20（场论的挠不变量）]
沿用注 3.25，由 $\EBord_{2|1}(\cX)$ 编码的超模空间可能相当微妙。特别地，带映入 $\cX$ 的映射的超欧几里得环面的模空间出乎意料地丰富。当 $\cX=X$ 为光滑流形时，[BE23a] 证明了该模空间上的函数决定 $H(M;\mathrm{MF})$ 中一个类的上闭链代表，即取值于模形式的通常上同调。超环面模空间的微妙特征在上闭链代表中编码全纯反常消去数据，从而恢复物理讨论 [GJFW21, GJF23] 所预期的信息。连同 Witten 对超对称 $\sigma$ 模型的局部构造（见例 4.3），这些反常消去数据可用于给出 Bunke 与 Naumann 的 $4k-1$ 维 string 流形挠不变量的一个量子场论解释 [BN14]。这些想法在 [BE23c] 中被置于 [ST11] 的框架之中，给出由 2 维超对称场论捕捉到的 TMF 挠的首批例子。
\end{example}
% [Tutor Note: 本段写 "For $\cX=X$ a smooth manifold"，与 §1 末尾同样的记号混用（前一个 $\cX$ 是叠）；照译作 $\cX=X$，已登记 ERRATA。]
% [Tutor Note: 上闭链所在的上同调原文写作 $H(M;\mathrm{MF})$，但本段通篇以 $X$ 记流形、$M$ 未定义；疑为 $H(X;\mathrm{MF})$ 之误，照排不改，待与 [BE23a] 原文核对（已记入回报）。]

% [原文 p.29]
\figomit{8}{全扩张的 $2|1$ 维欧几里得配边范畴中 2-态射的水平复合。条带 $S_{\ell,\tau,\bar\tau,\theta}$ 由入射弦与出射弦的长度、以及 $\HH^{2|1}\subset\R^{2|1}$ 中一个刻画条带角度与高度的参数所标定。文中把它们视为沿向上（时间）方向传播的开弦，并带有沿水平（空间）方向的粘合操作。把一条条带的边界粘合起来，就回到先前考虑过的柱面范畴。}

\begin{example}[4.21（次数扭与椭圆 Thom 类的对称性）]
与例 3.15 类似，定义 4.15 的次数扭可以构造为 $n$ 个手征自由费米子理论的量子化 [ST04, §2.6]，该理论由场与作用量泛函
\[
\psi\in\Gamma(\Sigma;\cS\otimes\R^n),
\qquad
S(\psi)=\frac{1}{2}\int_\Sigma\langle\psi,\bar\partial\psi\rangle
\]
定义，其中 $\Sigma$ 是带旋量丛 $\cS\to\Sigma$ 与手征 Dirac 算子 $\bar\partial$ 的 spin Riemann 面。该理论具有由映射 $\Sigma\to\mathrm{O}(n)$ 给出的经典（规范）对称性，作用方式为把 $\psi$ 在 $\R^n$ 中旋转。在量子理论中，这导致环路群 $L\mathrm{O}(n)$ 在状态 Hilbert 空间上的射影作用，它与 Clifford 代数 $\mathrm{Cl}(L\R^n)$ 的作用交织。这个环路群表示正是 $\widetilde{L\mathrm{Spin}}(n)$ 的级别 1（level one）表示的费米子 Fock 空间构造，其中 $\widetilde{L\mathrm{Spin}}(n)$ 是环路群 $L\mathrm{Spin}(n)$ 由 $p_1/2\in H^4(B\mathrm{Spin}(n);\Z)$ 分类的中心扩张 [PS86, §12]。这个手征共形场论决定一个具有 $\mathcal N=(0,1)$ 超对称的欧几里得场论。在猜想 1.1 之下，预期它恢复泛椭圆 Thom 类，见 [ST04, Theorem 1.0.3] 与 [DH11]。
\end{example}
% [Tutor Note: $p_1/2$ 取整性：$p_1$ 是 $B\mathrm{SO}(n)$ 的整类，拉回到 $B\mathrm{Spin}(n)$ 后被 2 整除，故 $p_1/2\in H^4(B\mathrm{Spin}(n);\Z)$，它正是该中心扩张的分类类。]
% [Tutor Note: 「泛椭圆 Thom 类」指与定义 4.15 的次数扭相配的 Thom 类，其在 string 定向下的取值为椭圆亏格；例 4.21 的作用是把这一几何对象实现为自由费米子场论，这正是 §4.5 中全扩张场论要恢复的对象。]

\subsection{最后的边疆：向下扩张}
全扩张的 $2|1$ 维欧几里得配边 2-范畴编码带角 $2|1$ 维超欧几里得流形的切割与粘合：由图 7 中的操作所编码的闭弦理论被细化为图 8 所描绘的开弦理论。预期是：向下扩张会把假设 4.18 中描述的取值于 $\KO((q))$ 与 $\mathrm{MF}^{\Z}$ 的不变量，细化为取值于 TMF 的场论不变量，即 (1) 中的上闭链映射。

这类全扩张场论的一些粗略特征已知，例如 [Moo03, 3.5] 与 [Seg01, §4]。然而仍有许多谜团：向下扩张不是唯一确定的，且尚无候选定义被详细写出。按惯例，这项工作可分为两个相互关联的部分：(i) 配边 2-范畴的几何；(ii) 目标 3-范畴的分析。$n$ 重 Segal 空间的语言提供了一个有前景的框架，可用来研究一般的全扩张配边范畴（例如见 [JFS17, GP20]），特别是全扩张的 $2|1$ 维欧几里得配边。目标的一个候选是共形网的对称幺半 3-范畴 [DH11, BDH15, BDH17, BDH19, BDH18]，不过预期需要一些技术性修改，因为这个 3-范畴并不去环（delooping）定义 1.5 与 [ST11] 中 1-扩张场论的 2-范畴 $\Alg$。

% [原文 p.30]
指导向下扩张中种种选择的一条途径，是更好地把握例子。或许最根本的例子是次数扭 (54) 及其与椭圆 Thom 类的预期联系，如例 4.21 所勾勒。它推广了 K 理论中已知的结构，见例 3.15。确实，一个实现如下形式的泛扭 Thom 同构的全扩张次数 $n$ 场论
\[
\QFT^{k}_{2|1}\left(\mathrm{pt}/\!/^{\nabla}\mathrm{O}(n)\right)\to\QFT^{n+k+\lambda}_{2|1}\left(\R^n/\!/^{\nabla}\mathrm{O}(n)\right)_c,
\qquad
E\mapsto E\boxtimes\Th_n
\]
将自动从 2 维超对称场论构造出一个 string 定向的可乘上同调理论。依照 Lurie 在 [Lur09b, §5.5] 中的建议，场论中的高阶等变结构随后可用来构造与 TMF 的比较，如 (1) 中那样。高阶等变性的证据来自如下预期：自由费米子场论的对称性是 string 2-群 [ST04, DH, SP11, Wal13]。

我们再加一个猜测。令 $\chi\QFT^{\bullet}_2\subset\QFT^{\bullet}_2$ 记次数为 $n$ 的理论中在 (55) 下的像为手征（例 4.2 意义下）者构成的子范畴。如推论 4.13 那样，预期这样的手征共形场论决定超对称的理论，从而给出一个包含
\begin{equation}
\chi\QFT^{\bullet}_2\xrightarrow{\ \text{群完备化？}\ }\QFT^{\bullet}_{2|1}.
\tag{56}
\end{equation}
与 (28) 类比，人们或可希望 (56) 实现手征场论的群完备化。确实，在 1 维，拓扑场论正是在重整化群流下不变的那些；而在 2 维，这些是手征共形理论。连同猜想 1.1，(56) 将把 TMF 实现为手征共形场论的群完备化。目前尚不清楚如何解读这一陈述，但很可能 (56) 的像中的刚性对象在猜想 1.1 中扮演重要角色，正如向量丛是 K 理论中的基本几何对象。
% [Tutor Note: (56) 的箭头标签原文为 "group completion?"，带问号，表示这是猜测而非已证的映射；译本照译作「群完备化？」。箭头因此不宜读作函子。]
% [Tutor Note: $\chi\QFT^{\bullet}_2$ 的上标用 $\bullet$，而正文说「次数为 $n$ 的理论」：$\bullet$ 表示对一切次数 $n$ 取并，逐次数的子范畴是 $\chi\QFT^{n}_2\subset\QFT^{n}_2$，与 p30 上文的 $\chi\QFT^n_2$ 记号一致。]
% [Tutor Note: 「1 维的拓扑场论恰是重整化群流下不变者」这一断言解释了为何 (56) 被视为群完备化：在 1 维，拓扑场论即完全可扩张到点的那部分；在 2 维，对应物是手征共形理论，而刚性对象（其对偶存在）扮演 K 理论中向量丛的角色。]
